\documentclass[11pt]{article}
\usepackage[margin=1in]{geometry}
\usepackage{amsmath,amssymb,amsthm,booktabs}
\usepackage[hidelinks]{hyperref}
\newtheorem{theorem}{Theorem}
\newtheorem{lemma}[theorem]{Lemma}
\newtheorem{corollary}[theorem]{Corollary}
\newtheorem{proposition}[theorem]{Proposition}
\newcommand{\C}{\mathcal C}
\newcommand{\bits}{\{0,1\}}
\newcommand{\wt}{\operatorname{wt}}
\title{Exact distances in marked-state information quotients\newline of the \v{C}ern\'{y} automata}
\author{ORION SyncLab research draft}
\date{28 September 2026}
\begin{document}
\maketitle
\begin{abstract}
For the \v{C}ern\'{y} automaton, retain subset cardinality and membership at
the consecutive marked states $0,n-1,\ldots,n-r+1$. Let $B_r(n)$ be the
distance from the full-set class to singleton classes in the existential
information quotient. We prove, for every $r\ge2$ and $n\ge r+2$,
\[
B_r(n)=3n+E(r)-2r-3,
\]
where $E(r)$ is the minimum external path length of a full binary tree with
$r$ leaves. A merger genealogy yields the lower bound through delayed
binary releases. An explicit reverse schedule achieves the bound with only
two hidden occupied states, with every input proved realizable. We derive
exact dyadic information increments and all fixed-codimension boundaries.
We also show that the same anchor count admits orders $n$, $n\log n$, and
$n^2$, by changing positions. The general formula is proved on paper;
the full optimal word language has a separate, sharp dyadic capacity
threshold, which we also determine.
Selected graph, potential, and timing lemmas are independently kernel-checked
in Lean. This is an abstraction theorem, not a proof of the \v{C}ern\'{y}
conjecture. Publication priority remains incompletely audited.
\end{abstract}

\section{Definitions and the exact quotient}
The automaton $\C_n$ has states $Q=\{0,\ldots,n-1\}$ and letters
$a(q)=q+1\pmod n$, $b(n-1)=0$, $b(q)=q$ otherwise. For an anchor set $A$,
let $\alpha_A(S)$ record $|S|$ and all memberships $q\in S$, $q\in A$,
for nonempty $S\subseteq Q$. A quotient edge $x\xrightarrow{c}y$ exists
if some $S$ has $\alpha_A(S)=x$ and $\alpha_A(Sc)=y$. Representatives may
change independently between successive edges. Write $B_A(n)$ for the
shortest quotient path from $\alpha_A(Q)$ to any singleton feature.
For $A_r=(0,n-1,\ldots,n-r+1)$ write $B_r(n)$.

Put $m=n-r$. A vertex is $(j,x)$, where $0\le j\le m$ counts occupied
hidden sites $1,\ldots,m$, $x\in\bits^r$ records the anchors in that order,
and $j+\wt(x)>0$. Writing $x=(x_0,\ldots,x_{r-1})$, its exact edges are
\begin{align}
(j,x)&\xrightarrow a
(j+x_0-t,(x_1,\ldots,x_{r-1},t)),
\quad t\le j\le m-1+t,\quad t\in\bits;\label{eq:a}\\
(j,x)&\xrightarrow b (j,(x_0\vee x_1,0,x_2,\ldots,x_{r-1})).\label{eq:b}
\end{align}
Indeed $t$ is occupancy of the exiting hidden site $m$. The other $m-1$
hidden sites may contain precisely $j-t$ particles, giving the displayed
necessary and sufficient condition. Letter $b$ acts only on the first two
anchors. Singleton goals satisfy $j+\wt(x)=1$.

\begin{proposition}[Orbit semantics]
For any finite DFA and any $A$, this quotient distance equals paid-letter
distance in the concrete power graph augmented with free permutations of
$Q\setminus A$ fixing $A$ pointwise.
\end{proposition}
\begin{proof}
Feature classes are exactly the orbits of that symmetric group: two sets
have identical observed parts and equal hidden cardinality precisely when
a hidden permutation takes one to the other. Augmented paths project to
quotient paths. Conversely, before each quotient edge, a free permutation
takes the current representative to its edge witness. Its letter realizes
that edge. Full-set and singleton endpoints are preserved. Commutation of
letters with permutations is unnecessary. The same statement holds for any
group whose orbits are the feature classes.
\end{proof}

\section{Main theorem and a short genealogy lower bound}
\begin{theorem}\label{thm:main}
For $r\ge2$, $n\ge r+2$, and $L=\lfloor\log_2(r-1)\rfloor$,
\[
B_r(n)=3n+Lr-2^{L+1}-3=3n+E(r)-2r-3.
\]
\end{theorem}

Lift any quotient path using orbit semantics, labelling each initial
particle and labelling every merged particle by the union of its ancestral
labels. An effective $b$ merges the occupied particles at $0,n-1$.
There are exactly $n-1$ effective letters on a path to a singleton, forming
a full binary genealogy. Count all other letters, including ineffective
$b$ letters, in a clock of total duration $T$. The path length is $n-1+T$.

\begin{lemma}[Timing constraints]\label{lem:timing}
Successive effective events differ by at least two clock units. A particle
created by merging at $0$ must wait at least $r$ clock units to participate
again at $n-1$, and at least $r+1$ to participate again at $0$.
\end{lemma}
\begin{proof}
After an effective $b$ the first two bits are $10$. One $a$ makes the first
bit zero; one ineffective $b$ leaves $10$. Neither permits another merge.
Contract the hidden block to $H$. A tracked particle moves at most one
step per clock letter around
\[
0\longrightarrow H\longrightarrow n-r+1\longrightarrow\cdots
\longrightarrow n-1\longrightarrow0.
\]
Hidden permutations stay at $H$, and effective events not involving that
particle leave it fixed. Reaching $n-1$ from $0$ needs $r$ steps. To merge
again at $0$ it cannot simply remain there: replenishing the emptied $n-1$
requires $a$, which moves it out of $0$. It must complete the cycle of
length $r+1$. An ineffective $b$ moving $n-1$ to $0$ is one clock step and
does not evade either bound.
\end{proof}

Read the genealogy backwards. Let $s_i$ be the clock of its $i$th split,
$1\le i\le M=n-1$. We have $s_1\ge0$ and $s_{i+1}\ge s_i+2$.
Each split releases its children for splitting no earlier than $s_i+r$
and $s_i+r+1$. These relaxed release times are ordered in parent order:
\[
0;\ s_1+r,s_1+r+1;\ s_2+r,s_2+r+1;\ldots .
\]
Every split consumes a distinct available genealogical particle. Even if
some children are never split, the $k$th consumption cannot precede the
$k$th relaxed release. Hence
\begin{equation}\label{eq:release}
s_k\ge s_{\lfloor k/2\rfloor}+r+(k\bmod2),\qquad k\ge2.
\end{equation}
In particular $s_{2i}\ge s_i+r$. Induction yields $s_{2^d}\ge dr$ for
$2^d\le M$. Spacing the remaining events gives
\[
T\ge s_M\ge dr+2(M-2^d).
\]
Take $d=L$, allowed since $2^L\le r-1\le M$. Adding $M=n-1$ proves
\begin{equation}\label{eq:lower}
B_r(n)\ge3n+Lr-2^{L+1}-3.
\end{equation}
This lower bound is valid even for $n=r+1$.

For completeness, the earliest relaxed clock recurrence
$t_1=0$, $t_k=\max(t_{k-1}+2,t_{\lfloor k/2\rfloor}+r+(k\bmod2))$
has the stronger exact solution
\[
t_k=2(k-1)+\max_{0\le d\le\lfloor\log_2k\rfloor}
\bigl(dr-2^{d+1}+2\bigr).
\]
To check it, subtract $2(k-1)$: only even indices create new maxima, giving
$f_k=\max(0,\max_{i\le k/2}(f_i+r-2i))$. A chain of $d$ choices has
indices at least $2^{d-1},\ldots,1$; the powers-of-two chain attains their
minimum sum $2^d-1$. This relaxed earliest schedule may exceed hidden
capacity two. We use a different schedule for the upper bound.

\section{A realizable capacity-two construction}
Set $K=r+1$, $C=Lr-2^{L+1}$, and
\[
F(i)=\min_{d\ge0}\bigl(2^{d+1}i-dr\bigr),\quad
s_i=C+F(i)\ (1\le i\le K),\quad A=s_K=C+2K.
\]
The minimum exists since the expressions eventually increase without
bound. Start at a singleton at anchor $0$ in the reverse quotient, with
hidden occupancy zero. Split it by reversing $b$ at clock zero. For every
integer $1\le t\le A$, reverse $a$ with input $u_t=1$ precisely when
$t\in\{s_2,\ldots,s_K\}$; otherwise use $u_t=0$. After every input one,
reverse $b$ to split at $0$.

The reverse rotation has exact input condition
\begin{equation}\label{eq:reverse}
u\le j\le1+u,
\end{equation}
and sends $(j,x)$ to $(j-u+x_{r-1},(u,x_0,\ldots,x_{r-2}))$.
Reverse splitting is legal exactly when the first two bits are $10$.
We prove all these requirements, including the upper bound in
\eqref{eq:reverse}; conservation alone would not suffice.

\begin{lemma}[Clock arithmetic]\label{lem:clock}
$s_1=0$, $s_{i+1}\ge s_i+2$, and $F(i)=2i$ when $2i\ge r$.
For every even $2p\le K$,
\[
s_{2p}-s_p=r,
\]
except when $r$ is odd and $2p=r+1$, where the difference is $r+1$.
\end{lemma}
\begin{proof}
Consecutive expressions in the defining minimum differ by $2^{d+1}i-r$.
The sign change gives $F(1)=2^{L+1}-Lr$, and gives $F(i)=2i$ when
$2i\ge r$. All affine expressions have slope at least two, so the same
spacing bound holds for their minimum. If $2p\le r$, a minimizing index
for $F(p)$ can be chosen at least one, since its $d=1$ term is no larger
than its $d=0$ term. Reindexing that minimum gives $F(2p)=F(p)+r$.
In the sole remaining case $2p=r+1$, both $F(p)=2p$ and $F(2p)=4p$,
giving the difference $r+1$.
\end{proof}

\begin{lemma}[Interlacing and realizability]\label{lem:buffer}
Every scheduled reverse edge is realizable at hidden capacity two.
\end{lemma}
\begin{proof}
A split at $s_p$ produces particles at $0,n-1$. Subsequent reverse
rotations release them into the hidden block after $r$ and $r-1$ steps,
respectively; their first legal input times are therefore $s_p+r+1$ and
$s_p+r$. Later splits fill an empty site $n-1$ and do not change these
trajectories. Include a dummy root ready time $q_1=0$. The ordered input
ready times are
\[
q_{2p}=s_p+r,\qquad q_{2p+1}=s_p+r+1.
\]
Spacing makes this list strictly increasing. Lemma~\ref{lem:clock} gives
$q_k\le s_k$ for $2\le k\le K$: even indices are immediate and odd
indices follow by spacing. It also gives $s_{k-1}\le q_k$ for
$2\le k\le K$. For even $k=2p$, use
$s_{2p-1}\le s_{2p}-2\le s_p+r-1$; for odd $k=2p+1$, use
$s_{2p}\le s_p+r+1$.

The endpoint index must be checked separately. If $r=2p$, then
$q_{K+1}=s_{p+1}+r=A$ and $q_{K+2}=s_{p+1}+r+1=A+1$.
If $r=2p+1$, then $q_{K+1}=s_{p+1}+r+1=A$ and
$q_{K+2}=s_{p+2}+r=A+1$. All involved $F$ values lie in the $2i$ regime.
Thus $s_K\le q_{K+1}$ holds as equality.

Let $R(t)=\#\{q_i\le t\}$ and $N(t)=\#\{s_i\le t\}$. Interlacing gives
\[
N(t)\le R(t)\le N(t)+1\qquad(0\le t\le A).
\]
Before rotating at time $t\ge1$, the hidden count is
$j_{t-1}=R(t)-N(t-1)$: exits occur one unit before input readiness and
previous inputs are exactly previous non-root splits. The dummy root and
its split cancel in this count. If $u_t=1$, then
$N(t)=N(t-1)+1$, giving $1\le j_{t-1}\le2$. If $u_t=0$, the counts
agree, giving $0\le j_{t-1}\le1$. This is precisely
\eqref{eq:reverse}. These identities validate the process inductively.
Finally, before each non-root split the preceding input was zero, since
split clocks differ by at least two. The first two bits after its rotation
are consequently $10$, so the reverse $b$ is legal too.
\end{proof}

\begin{lemma}[Full endpoint]
The construction ends at the all-one anchor window with hidden count two.
\end{lemma}
\begin{proof}
For $i>r/2$, $s_i=C+2i$. A split fills consecutive stream positions
$s_i,s_i-1$. There are at least $\lceil r/2\rceil+1$ terminal splits in
this regime, so the last $r$ positions $A,A-1,\ldots,A-r+1$ are all filled.
Each split adds one particle, and there are $K=r+1$ splits from one initial
particle. Thus $r+2$ particles are present, of which $r$ occupy anchors;
exactly two are hidden. Equivalently the strict release ordering and
$q_{K+2}=A+1$ give hidden count $R(A+1)-N(A)=2$.
\end{proof}

The reverse path has $A$ rotations and $K$ splits, hence length
$A+K=C+3K=(L+3)r-2^{L+1}+3$. This meets \eqref{eq:lower} at $n=r+2$.
For arbitrary $m\ge2$, prepend $m-2$ forward $baa$ cycles, reducing the
full window's hidden count from $m$ to two. Formula~\eqref{eq:a} shows a
path valid at capacity two is valid at every larger capacity. Adding
$3(m-2)$ proves Theorem~\ref{thm:main}.

\section{Binary trees and the information hierarchy}
\begin{lemma}[Classical external path length]
For $r\ge2$, $h=\lceil\log_2r\rceil$,
$E(r)=(h+1)r-2^h$.
\end{lemma}
\begin{proof}
If an optimal full binary tree has a shallowest leaf at depth $d$ and
deepest leaves at depth $D\ge d+2$, choose a deepest sibling pair. Replace
that pair by its parent leaf and split the shallow leaf. The leaf count is
unchanged and the external path length changes by $d-D+1<0$, a
contradiction. Thus depths differ by at most one. Kraft equality gives
$2^h-r$ leaves at depth $h-1$ and $2r-2^h$ at depth $h$, with the usual
all-equal case when $r$ is a power of two. Such a tree exists: split
$r-2^{h-1}$ leaves of the complete tree at depth $h-1$. Summing depths
gives the formula. This is the equal-weight prefix-code optimization
problem, not a new coding theorem.
\end{proof}

\begin{corollary}
For $2^{h-1}<r\le2^h$ and $n\ge r+2$,
\[
B_r(n)=3n+(h-1)r-2^h-3.
\]
For $2\le r\le n-3$, $B_{r+1}(n)-B_r(n)=\lfloor\log_2r\rfloor$.
For fixed $k\ge2$, $n\ge k+2$, setting $h=\lceil\log_2(n-k)\rceil$ gives
\[
B_{n-k}(n)=n+2k-3+E(n-k)=(h+2)n-(h-1)k-2^h-3.
\]
In particular $B_{n-2}(n)=n+1+E(n-2)=\Theta(n\log n)$.
\end{corollary}
These follow directly by substitution; $E(r+1)-E(r)=\lfloor\log_2r\rfloor+2$,
including the powers-of-two boundary. For $r=n-1$ the feature is injective,
so the classical threshold theorem \cite{volkov} gives
$B_{n-1}(n)=(n-1)^2$. The last increment is therefore
$n^2-3n-E(n-2)=\Theta(n^2)$ for $n\ge4$.
The dependence on hidden capacity stabilizes exactly at two:
$B_r(n+1)-B_r(n)=3$ for every $n\ge r+2$.

Writing $x=r/2^h\in(1/2,1]$ also gives
$E(r)=r\log_2r+r(1-1/x-\log_2x)$.
The parenthesized nonnegative term attains its maximum
$1-1/\ln2-\log_2(\ln2)$ at $x=\ln2$, by differentiation. Thus the
linear error contains an explicitly bounded dyadic oscillation.

\section{Optimal genealogies and the full optimal word language}
\begin{proposition}[Greatest relaxed schedule]
For $r,K\ge1$, consider $s_1\ge0$, $s_{i+1}\ge s_i+2$, and
$s_{2i}\ge s_i+r$ whenever $2i\le K$. The minimum final clock is
\[
T(r,K)=2K+\max_{0\le d\le\lfloor\log_2K\rfloor}(dr-2^{d+1}).
\]
For every deadline $A\ge T(r,K)$ the unique componentwise greatest
feasible vector is
\[
g_i=A-2K+\min_{2^di\le K}(2^{d+1}i-dr).
\]
\end{proposition}
\begin{proof}
Repeated doubling and final spacing give
$s_i+dr\le s_{2^di}\le A-2(K-2^di)$, so $s_i\le g_i$.
Write the minimum as $G(i)$. Its domain shrinks with $i$, and every term
has slope at least two, proving spacing. Reindexing gives
$G(2i)=r+\min_{e\ge1,2^ei\le K}(2^{e+1}i-er)\ge r+G(i)$.
Finally $g_K=A$ and $g_1=A-T(r,K)\ge0$, proving feasibility and
attainment. A smaller deadline makes $g_1<0$, impossible.
For $K\ge r\ge2$, this specializes to $T=2K+C$, $g_i=C+F(i)$.
The minimum has an unrestricted minimizer with $2^di\le r-1\le K$,
or $d=0$ in the $2i\ge r$ regime.
\end{proof}

\begin{theorem}[Rigidity and no wasted letters]\label{thm:rigidity}
Put $p=2^L$. Every optimal quotient path, $n\ge r+2$, has reverse clocks
\[
s_{2^d}=dr\ (0\le d\le L),\qquad s_i=C+2i\ (p\le i\le n-1).
\]
All its $b$ letters are effective. It starts with $(baa)^{n-1-p}b$.
For $r=2,3,4$ its optimal letter word is unique:
\[
(baa)^{n-2}b;\qquad (baa)^{n-3}ba^3b;\qquad(baa)^{n-3}ba^4b,
\]
respectively, although representatives need not be unique.
\end{theorem}
\begin{proof}
Equality in $T\ge s_{n-1}\ge s_p+2(n-1-p)\ge Lr+2(n-1-p)$ forces
the dense tail and every doubling step of the spine. It also forces
$s_1=0$ and $s_{n-1}=T$, so there are no clock letters beyond the event
interval. Label the strictly ordered relaxed releases by root entry one,
the $n-1$ child of event $i$ as entry $2i$, and its $0$ child as entry
$2i+1$. If $s_k=q_k$, no release entry beyond $k$ can have supplied any
of the first $k$ consumed particles. Thus exactly entries $1,\ldots,k$
were consumed, and entry $k$ must be consumed by event $k$ itself, since
earlier clocks are less than its readiness bound. For $k=2^d$ we have
$s_k=s_{k/2}+r=q_k$, so the dyadic event is the actual $n-1$ child of
its predecessor on the spine. Its waiting time $r$ is tight. Traversing
the $r$ directed window steps in $r$ clock letters uses only $a$: an
ineffective $b$ either stays fixed or returns $n-1$ to $0$. These intervals
cover clocks zero through $s_p$.
Every remaining gap is two clock letters. From post-merge bits $10$,
neither $ab$, $ba$, nor $bb$ with ineffective $b$ permits another merge;
only $aa$ can. Thus every clock letter is $a$. The dense tail becomes
the asserted forward prefix. For $r=2,3,4$, all event indices are forced,
so the listed words are the only possibilities.
\end{proof}

Set $K_0=r+1$, $A_0=C+2K_0$. Let $W_r$ consist of letter words whose
reverse split clocks satisfy $s_1=0$, $s_{K_0}=A_0$, spacing two, and
even releases $s_{2i}\ge s_i+r$. Odd releases follow by spacing.
For hidden capacity $m\ge2$, every optimum starts with $(baa)^{m-2}$.
Strip that prefix and call its reduced language $R_r(m)$.

\begin{theorem}[Sharp language stabilization]\label{thm:language}
The languages increase as
$R_r(2)\subseteq R_r(3)\subseteq\cdots\subseteq W_r$.
The least capacity $M_r$ from which they equal $W_r$ is
\[
M_2=2,\qquad M_r=1+2^{L-1}\quad(r\ge3).
\]
\end{theorem}
\begin{proof}
Necessity follows from timing and rigidity: the first $r+1$ reverse
events of every optimum have deadline $A_0$. Padding gives the nesting.
For $r\ge3$, write $p=2^L$, $M=p/2+1$. Take any $W_r$ vector and
append $m-2$ dense events, giving $K=r+m-1$. The extended vector still
satisfies even releases: for a new child $2i>r+1$, its parent is at most
$C+F(i)=C+2i$ and its clock is $C+4i$. Run input one at splits and zero
elsewhere. With the ordered releases $q$ and counts $R,N$, $N(t)\le R(t)$
and hidden count before rotation is $j=R(t)-N(t-1)$. This proves the
input lower bound; spacing proves the split bits $10$.

Here are the upper bounds, retaining the crucial one-unit difference
between zero and one input. Put $a=s_p=Lr$. If $t<a$, then
$i=N(t-r)\le p/2-1$, $R(t)\le1+2i$, $N(t-1)\ge i$, so $j\le p/2$.
Negative arguments of $N$ are interpreted as zero. At $t=a$, $j=1$.
For $t=a+\delta$, $1\le\delta<r$, spacing from event $p/2$ gives
\[
i=N(t-r)\le\min(p-1,p/2+\lfloor\delta/2\rfloor),\quad
N(t-1)=p+\lfloor(\delta-1)/2\rfloor.
\]
At even $\delta=2z$, a split, this yields
$j\le2\min(p-1,p/2+z)-p-z+2\le p/2+1=M$.
At odd $\delta=2z+1$ it yields
$j\le2\min(p-1,p/2+z)-p-z+1\le p/2=M-1$.
Both inequalities follow by splitting at $z=p/2$.
For $\delta\ge r$, the dense parents release one particle per clock:
$R(t)=2p+\delta-r$. At $\delta=2z$ or $2z+1$, $j=p-r+z+1$.
A split has $z\le K-p$, hence $j\le m$; a nonsplit has
$z\le K-p-1$, hence $j\le m-1$. Thus every input is realizable for
$m\ge M$. The last $K-p+1$ dense splits fill
$2(r+m-p)\ge r$ stream positions; the endpoint has full window and
$K+1=n$ particles, so hidden count $m$. This proves $R_r(m)=W_r$.

For the matching obstruction, let $r\ge5$, $m=M-1=p/2$, and choose
the earliest relaxed recurrence, extended to $K=r+m-1$. All its inputs
are forced: split input is one, the immediately preceding input is zero,
and in any longer waiting gap the next split waits for its next release,
so hidden occupancy at all other inputs is zero. In particular no
unpaired input can avert overflow. The earliest formula gives
$s_{p-1}=(L-1)r+p-2$. At the nonsplit clock
$t^*=Lr+p-1$, immediately before event $p+p/2$ (which exists since
$r\ge p+1$),
\[
q_{2p-1}=t^*<q_{2p},\quad R(t^*)=2p-1,\quad
N(t^*-1)=3p/2-1,\quad j=p/2=m.
\]
Its forced zero input requires $j\le m-1$, impossible. Hence this $W_r$
word is missing at $M-1$. For $r=3,4$ the upper threshold is already
the minimum allowed capacity two; $r=2$ follows from its unique word.
\end{proof}

The first strict enlargement is $r=5$: the word
$(baa)^2ba^3ba^2ba^5b$ of length 20 is impossible at capacity two, but
prepending $baa$ gives an optimal capacity-three word of length 23.
Its reverse clocks are $0,5,7,10,12,14$; before clock 13 the hidden count
is two and the split at 14 forces input zero, violating capacity two.
At capacity three, append a split at 16 and all inputs are realizable.
Thus value stabilization and optimal-language stabilization are distinct.

\begin{proposition}[End-of-regime language bijections]
For $L\ge2$, $p=2^L$, the languages $W_{2p-2},W_{2p-1},W_{2p}$ have
explicit bijections and common size $C_L\ge2^{p/2-1}$.
\end{proposition}
\begin{proof}
For $i=2^d+a<p$, normalize clocks as
$s_i=dr+2a+z_i$; set $z_i=0$ at every power of two, and at indices
$i\ge p$. The common conditions on the nonnegative integer array are
monotonicity within each binary-index level and
$z_i\le z_{2i}+2a$. All children $2i$ of indices $i<p$ exist for the
three indicated $r$. Iteration gives
$z_i\le2a(2^{L-d}-1)$, so the remaining boundary spacing condition
$z_{2^{d+1}-1}\le r-2^{d+1}$ is automatic:
\[
2(2^d-1)(2^{L-d}-1)
=2p-2^{d+1}-2^{L-d+1}+2\le2p-2-2^{d+1}.
\]
Tail constraints are automatic as well. Thus this same finite array set
parametrizes each language bijectively. For the lower count set upper
levels to zero and choose arbitrary bits $b_1,\ldots,b_{p/2-1}$ on the
bottom level through $z_{p/2+a}=\sum_{j\le a}b_j$. These satisfy all
conditions and yield the stated number of distinct arrays. Complete finite
enumeration gives $C_2=3,C_3=279$; no general numerical count is claimed.
\end{proof}

\section{Anchor placement and a transferable lower-bound principle}
\begin{proposition}[Observed interval charge]
If an anchor layout contains an interior interval $I\subseteq\{1,\ldots,n-2\}$,
then $B_A(n)\ge n-1+|I|$.
\end{proposition}
\begin{proof}
Truncate any synchronizing quotient path after its last effective $b$.
Its goal is the singleton at $0$. The occupied count in $I$ is feature
determined, is fixed by $b$, and decreases by at most one per $a$ (only the
last interval site can exit). It starts at $|I|$ and ends at zero, requiring
at least $|I|$ rotations in addition to $n-1$ effective letters. Truncation
cannot increase length. For the empty interval this is the cardinality bound.
\end{proof}

\begin{theorem}[An exact linear family]\label{thm:linear}
For $2\le k\le n-1$, hide $H=\{1,\ldots,k-1,n-1\}$ and observe its
complement $A=\{0,k,\ldots,n-2\}$. Then
\[
B_A(n)=2n-k-2=n+|A|-2.
\]
\end{theorem}
\begin{proof}
The observed interval $I=\{k,\ldots,n-2\}$ has size $n-k-1$, giving the
lower bound. A matching abstract word is $b^{k-1}(ab)^{n-k-1}b$.
The first $k-1$ effective letters leave all anchors full and hidden count
one: choose $n-1$ among the hidden occupants of each witness. In each
following rotation choose that unique hidden occupant at $n-1$.
Anchor $0$ stays occupied, its particle enters hidden site $1$, and the
occupied suffix of $I$ shifts right, losing its final particle into hidden
$n-1$. Hidden count becomes two, while the observed interval count drops
by one. For its ensuing effective $b$ choose hidden sites $\{1,n-1\}$;
hidden count returns to one and anchors are unchanged. After $n-k-1$ pairs
the interval is empty. The final effective $b$ leaves singleton $0$.
All witnesses exist because $k\ge2$. The length is $2n-k-2$.
\end{proof}

\begin{proposition}[One observed state]
For $n\ge2$, $B_{\{0\}}(n)=n-1$ and $B_{\{n-1\}}(n)=2n-3$.
For $n\ge3$ and $1\le q\le n-2$, $B_{\{q\}}(n)=n$.
\end{proposition}
\begin{proof}
Cardinality needs at least $n-1$ effective letters. When only $0$ is
marked, every effective $b$ edge after the first is witnessed by a subset
of the required size containing both endpoints, giving $n-1$.
When only $n-1$ is marked with bit $v$, the potential
$P(k,v)=2k-2-v$ for $k\ge2$, and zero for $k=1$, drops by at most one
per edge. Rotation changes $v$ by at most one; an effective $b$ sends
$(k,1)$ to $(k-1,0)$; ineffective cases also satisfy the inequality.
Its full value is $2n-3$. The matching path is $b(ab)^{n-2}$: each rotation
edge from $(k,0)$ to $(k,1)$ has a witness containing $n-2$ and excluding
$n-1$, and each $b$ edge has an endpoint-pair witness.
For interior $q$, a path of length $n-1$ would consist entirely of
effective $b$ letters, preserving the marked bit one. Its last two-set
would need to be $\{0,n-1\}$, which has marked bit zero, impossible.
For the upper bound, use $n-2$ effective $b$ edges down to cardinality two
with marked bit one, witnessed by sets containing $\{q,0,n-1\}$ before
each decrease. A two-set containing $q$ and excluding $q-1$ witnesses a
rotation to marked bit zero; a final endpoint-pair witness merges it.
The length is $n$.
\end{proof}

\begin{proposition}[Feature-resolved monotone simulation]\label{prop:sim}
Let $D$ be a synchronizing DFA and $D_D(U)$ its subset reset distance.
Suppose a nonempty-subset map $V(S)$ is determined by the source feature,
sends source goals to distance-zero subsets, and each concrete source
edge $S\to S'$ admits a subset $U$ reached from $V(S)$ by zero or one
$D$-letter with $U\subseteq V(S')$. Then $D_D(V(S))$ is a quotient
potential, and its initial value is a lower bound on quotient distance.
\end{proposition}
\begin{proof}
Subset reset distance is monotone under inclusion: any word resetting a
larger subset resets every smaller nonempty subset. Thus
$D_D(V(S))\le1+D_D(U)\le1+D_D(V(S'))$. Feature dependence transfers this
inequality to every existential edge, and telescoping proves the claim.
\end{proof}

\begin{theorem}[Position changes the asymptotic order]
If $g\ge2$ divides $n$ and all hidden states lie in one residue class
modulo $g$, then $B_A(n)\ge(g-1)^2$.
Consequently, for even $n$ and hidden states $\{q,q+n/2\}$,
$B_A(n)=\Theta(n^2)$. With the same number $n-2$ of anchors but hidden
states $\{1,2\}$, the value is $\Theta(n\log n)$; with hidden states
$\{1,n-1\}$ it is exactly $2n-4$.
\end{theorem}
\begin{proof}
Let $V(S)=\{q\bmod g:q\in S\}$. The feature determines $V(S)$:
every residue except the hidden residue is entirely observed, and the
hidden residue's occupancy is determined by its observed bits together
with whether hidden occupancy is positive. Projection commutes with $a$.
If $n-1\notin S$, the concrete $b$ is a loop and choose $U=V(S)$.
If $n-1\in S$, take $U=V(S)b_{\C_g}$; then $U\subseteq V(Sb)$, because
other representatives of residue $g-1$ may remain and the new zero is
present. Proposition~\ref{prop:sim} gives the classical full distance
$D_{\C_g}(Q_g)=(g-1)^2$ as a lower bound. For antipodal hidden states
choose $g=n/2$; the concrete reset word gives upper bound $(n-1)^2$.
The adjacent hidden pair is precisely $A_{n-2}$, covered by the main
theorem; the linear family is Theorem~\ref{thm:linear}. These are statements
along even $n$ tending to infinity.
\end{proof}
More generally, fix any $k\ge2$. Along $n\ge2k$ divisible by $k$, let $g=n/k$
and hide the progression $\{q,q+g,\ldots,q+(k-1)g\}$, $0\le q<g$.
Then $B_A(n)\ge(n/k-1)^2$ and $B_A(n)=\Theta(n^2)$. At the same count
$n-k$, hiding $\{1,\ldots,k\}$ gives $\Theta(n\log n)$, while hiding
$\{1,\ldots,k-1,n-1\}$ gives the exact linear value $2n-k-2$.

\section{Verification, limitations, and related work}
\paragraph{Mathematical status.}
Theorem~\ref{thm:main} and the layout theorems are \textbf{PROVED} by the
arguments above. No finite experiment is used as an induction hypothesis.
The complete formula is not labelled \textbf{FORMALLY VERIFIED}.
\paragraph{Independent computations.}
The search-free clock construction was checked by a concrete-set image
verifier at 111 values of $r$, including $2,3$, and numerous dyadic
boundaries through $1025$. It enumerates hidden-site choices and does not
use the quotient edge formula to validate witnesses. A different backward
latest-schedule recurrence agrees with the direct minimum formula.
Release interlacing arithmetic was checked at boundary values through
$r=65537$. In 1750 sampled complete abstract paths with random prefixes,
including 18446 ineffective $b$ letters, a labelled-particle checker found
no violation of the timing inequalities. Complete quotient BFS at
capacity two for $r=18,\ldots,23$ gave lengths $97,104,111,118,125,132$.
An independent concrete-subset enumeration checked all single-anchor
positions for $2\le n\le12$ and all hidden pairs for $4\le n\le12$.
The exact endpoint family was additionally checked by complete concrete
quotient BFS at 55 parameter points through $n=12$, and by explicit
set-image witnesses through $n=1000$.
These statements are \textbf{COMPUTATIONALLY VERIFIED} finite evidence.
For optimal-language claims, all optimal-path edges were explored in
66 small models: 3579 vertices and 3659 edges obeyed the forced-prefix
and no-wasted-$b$ restrictions. Complete enumeration of all 1352 relaxed
clock vectors for $r=2,\ldots,16$ at all relevant capacities matches the
sharp threshold. An additional 2700 random clock/capacity witnesses and
135 concrete-set reverse checks through $r=1025$ verified its upper
construction; the universal earliest obstruction occurred at exactly
the predicted clock for every tested $r\ge5$.

\paragraph{Kernel boundary.}
Lean 4.19.0 checks existential edge and arbitrary-walk potential theorems,
the monotone-simulation transfer, finite dyadic release/spacing arithmetic,
capacity-two input arithmetic, and two/three-anchor symbolic potential
inequalities, optimal-tail rigidity, and sharp dyadic buffer arithmetic.
It also checks the $3n-4$ lower bound for every walk in the
explicit three-anchor count/window graph, and concrete optimal subset
paths for $\C_3,\C_4$. No proof holes, added axioms, or
\texttt{native\_decide} are used. The general genealogy lifting/timing
bridge, closed-form clock interlacing, and window endpoint remain paper
proofs. No formal proof of the full main equality is claimed.

\paragraph{Reproduction.}
From the repository root, run
\begin{verbatim}
python post_baseline/check_uniform.py
python post_baseline/layout_experiments.py
python post_baseline/endpoint_layout.py
python post_baseline/check_normal_form.py
python post_baseline/check_language_stabilization.py
python post_baseline/check_sharp_language.py
python post_baseline/check_language_counts.py
python post_baseline/uniform_construction.py --r 33 --m 2
./post_baseline/lean_sprint/check.ps1 -LeanExe lean
\end{verbatim}
The Python theorem checkers use only the standard library. Optional SMT
explorations require Z3 but are unnecessary for the construction or proof.
The checked Lean source, individual compiler exit codes and axiom logs,
finite BFS data, failed-route diagnoses, and detailed proof companions are
retained in \texttt{post\_baseline/}. Frozen baseline files are unchanged.

\paragraph{Prior work and novelty.}
The quadratic threshold of $\C_n$ is classical \cite{volkov}, and exact
power-automaton search is established \cite{search}. Concrete subset
synchronization and compression are different questions surveyed in
\cite{openproblems}. Prefix-code decoders already have a substantial
synchronization theory \cite{huffman}; a general claim that binary trees
and synchronization are newly connected would be incorrect. Delayed
genealogies also resemble lopsided-tree broadcast models \cite{lopsided}.
The orbit semantics and simulation potential are standard arguments;
permutation-group synchronization is established \cite{groups}.
A bounded literature search located no matching all-parameter theorem for
this particular information quotient, capacity-two realization, or
same-count layout order separation. Their novelty is \emph{plausible but
not fully confirmed}, including the sharp optimal-language threshold.
In particular, priority for the isolated relaxed
event recurrence remains open pending a deeper scheduling-literature
comparison. This draft makes no first-discovery claim.

\paragraph{Open mathematical and formal questions.}
Determine exact values for general anchor layouts, especially antipodal
hidden pairs; characterize layouts with linear, $n\log n$, or quadratic
information distances; extend the capacity scheduling theorem to more
general feature partitions and automaton families; complete the Lean
genealogy and construction interfaces. The main contiguous-anchor formula
and sharp optimal-language threshold
have no remaining paper-proof gap, but these generalizations and full
formal verification are \textbf{OPEN}.

\begin{thebibliography}{9}
\bibitem{volkov} M. V. Volkov.
Synchronization of finite automata.
\emph{Russian Math. Surveys} 77(5), 819--891, 2022.
\href{https://www.mathnet.ru/eng/rm10005}{doi:10.4213/rm10005e}.
\bibitem{openproblems} M. Szyku{\l}a.
Synchronizing Automata: Open Problems, 2026.
\href{https://arxiv.org/abs/2608.24245}{arXiv:2608.24245}.
\bibitem{search} A. Kisielewicz, J. Kowalski, M. Szyku{\l}a.
Computing the shortest reset words of synchronizing automata.
\emph{J. Combin. Optim.} 29, 88--124, 2015.
\href{https://doi.org/10.1007/s10878-013-9682-0}{doi:10.1007/s10878-013-9682-0}.
\bibitem{huffman} M. V. Berlinkov, M. Szyku{\l}a.
Algebraic synchronization criterion and computing reset words.
\emph{Information Sciences} 369, 718--730, 2016.
\href{https://arxiv.org/abs/1412.8363}{arXiv:1412.8363}.
\bibitem{lopsided} M. Golin, A. Schuster.
Optimal point-to-point broadcast algorithms via lopsided trees.
\emph{Discrete Applied Math.} 93(2--3), 233--263, 1999.
\href{https://doi.org/10.1016/S0166-218X(99)00107-9}{doi:10.1016/S0166-218X(99)00107-9}.
\bibitem{groups} J. Ara{\'u}jo, P. J. Cameron, B. Steinberg.
Between primitive and 2-transitive: Synchronization and its friends.
\emph{EMS Surv. Math. Sci.} 4, 101--184, 2017.
\href{https://doi.org/10.4171/EMSS/4-2-1}{doi:10.4171/EMSS/4-2-1}.
\end{thebibliography}
\end{document}
